% This document was generated by an LLM.
% Content fragment; compile intro-analysis-self-check.tex.
% Coverage: Hunter (2012), Chapter 6, printed pp. 73--92.
% Source: https://www.math.ucdavis.edu/~hunter/m125a/intro_analysis.pdf
% Corrections and clarifications:
% - Treat the center separately when a radius or a ratio is zero.
% - In Example 6.12, the coefficient ratio is (n+1)/n, and the
%   endpoint series at x=0 is minus the harmonic series.
% - In Example 6.13, the nonzero coefficients have alternating signs:
%   a_(2^k)=(-1)^k, with all other coefficients, including a_0, zero.
% - For the converse radius argument in Theorem 6.14, use convergence
%   at a larger distance to obtain absolute convergence at smaller ones;
%   distinguish signed convergence from absolute convergence.
% - Repeated differentiation in Theorem 6.16 uses Theorem 6.15 and
%   retains the coefficient a_n in the formula for each derivative.
% - State Corollary 6.17 using equality near the common center c,
%   as in its proof, rather than the printed neighborhood of 0.
% - In Proposition 6.20 use (f/E)'=(f'E-fE')/E^2.
% - Only the zero analytic function on R can have compact support.

\section{Power Series}

\subsection{Introduction}
\textit{Source: \S6.1, pp.~73--74.}

\begin{enumerate}
  % Definition 6.1 and introductory discussion.
  \item Define a power series centered at \(c\), its coefficients,
    and its partial sums. When is it a polynomial?
    Why does writing down a formal series not yet define a function
    at every real input?

  % Examples and change of variables in Section 6.1.
  \item Give a power series centered at zero and one centered at one,
    identifying their coefficients. Explain how a change of variable
    reduces a convergence question at a general center to one at zero.
\end{enumerate}

\subsection{Radius of convergence}
\textit{Source: \S6.2, pp.~74--76.}

\begin{enumerate}[resume]
  % Theorem 6.2.
  \item State the convergence theorem for power series (Theorem 6.2).
    Specify where convergence is absolute, where divergence occurs,
    where uniform convergence is guaranteed, and where the sum is continuous.
    What does the theorem leave undecided at a finite positive radius?

  % First step in the proof of Theorem 6.2.
  \item Suppose \(\sum_{n=0}^{\infty}a_nx_0^n\) converges at
    \(x_0\ne0\). Prove absolute convergence for every \(|x|<|x_0|\)
    by bounding the terms at \(x_0\) and comparing with a geometric series.

  % Construction of R in the proof of Theorem 6.2.
  \item Define the radius at center zero as a supremum of distances
    where the series converges. Explain why this gives the asserted
    convergence and divergence regions. Treat radii zero and infinity
    explicitly, including convergence at the center.

  % Uniform convergence step in Theorem 6.2.
  \item For \(0\leq\rho<R\), prove uniform convergence on
    \(|x-c|\leq\rho\) by choosing \(\rho<\sigma<R\) and applying
    the \(M\)-test. How does continuity of the sum follow without
    assuming uniform convergence on the entire open convergence interval?

  % Definition 6.3 and endpoint discussion.
  \item Define radius of convergence and interval of convergence.
    For \(0<R<\infty\), list the possible endpoint patterns.
    What are the convergence sets when \(R=0\) and \(R=\infty\)?

  % Theorem 6.4.
  \item State the coefficient-ratio formula for a radius of convergence
    (Theorem 6.4), including the eventual nonzero-coefficient and
    limit assumptions. Give the absolute-value formula and its
    interpretation when the ratio tends to zero or infinity.

  % Proof of Theorem 6.4.
  \item Derive Theorem 6.4 from the numerical ratio test.
    How does the ratio of successive series terms depend on
    \(|x-c|\)? Why must the center be handled separately?

  % Theorem 6.5.
  \item State Hadamard's formula for the radius in terms of
    \(\limsup_{n\to\infty}|a_n|^{1/n}\) (Theorem 6.5), with
    its zero and infinite cases. Why does this formula apply even
    when the coefficient-ratio formula is unavailable?

  % Proof of Theorem 6.5.
  \item Derive Hadamard's formula from the root test.
    Explain why neither the root-test threshold nor the ratio-test
    threshold alone determines convergence at an endpoint.
\end{enumerate}

\Needspace{8\baselineskip}
\subsection{Examples of power series}
\textit{Source: \S6.3, pp.~76--79.}

\begin{enumerate}[resume]
  % Example 6.6.
  \item Find the radius, interval, and sum of
    \(\sum_{n=0}^{\infty}x^n\). Compare the convergence set with
    the domain of the rational function that gives the sum.
    Is convergence uniform on the entire convergence interval?

  % Example 6.7.
  \item Determine the radius and interval of convergence of
    \(\sum_{n=1}^{\infty}x^n/n\), testing both endpoints separately
    and distinguishing absolute from conditional convergence.
    Does the series converge uniformly on \((-1,1)\)?

  % Example 6.8.
  \item Find the radius and interval of convergence of
    \(\sum_{n=0}^{\infty}x^n/n!\). What familiar function will
    its sum define? What uniform convergence on bounded intervals
    follows from the radius theorem?

  % Examples 6.9--6.10.
  \item Find the radii and intervals of convergence of
    \[
      \sum_{n=0}^{\infty}\frac{(-1)^n x^{2n}}{(2n)!},
      \qquad
      \sum_{n=0}^{\infty}\frac{(-1)^n x^{2n+1}}{(2n+1)!}.
    \]
    Identify their sums. How can a ratio test on the displayed
    terms be used even though the full sequences of power-series
    coefficients contain zeros?

  % Example 6.11.
  \item Find the radius and convergence set of
    \(\sum_{n=0}^{\infty}n!x^n\).
    For a fixed \(x\ne0\), explain why the terms eventually grow
    in magnitude rather than tending to zero.

  % Example 6.12; correct the coefficient ratio and endpoint sign.
  \item Determine the center, radius, and interval of convergence of
    \[
      \sum_{n=1}^{\infty}\frac{(-1)^{n+1}}{n}(x-1)^n.
    \]
    Compute the coefficient ratio carefully and test each endpoint,
    including the signs of its terms.

  % Example 6.13; use the correct signed coefficient sequence.
  \item Write the lacunary series
    \(\sum_{k=0}^{\infty}(-1)^k x^{2^k}\) in the form
    \(\sum_{n=0}^{\infty}a_nx^n\) by specifying every coefficient.
    Find its radius and interval of convergence using comparison
    and then Hadamard's formula. Why is the usual coefficient-ratio
    formula unsuitable?

  % Definition and unproved boundary-limit observation in Example 6.13.
  \item What is meant by a lacunary power series?
    For the preceding example, distinguish convergence of the series
    at \(x=1\) from existence of the limit of its sum as \(x\to1^-\).
    What behavior of the sum near one is reported in the notes?
\end{enumerate}

\subsection{Differentiation of power series}
\textit{Source: \S6.4, pp.~79--82.}

\begin{enumerate}[resume]
  % Theorem 6.14.
  \item State the radius-preservation theorem for the formal
    term-by-term derivative of a power series (Theorem 6.14).
    Does it assert that the two series have the same endpoint behavior?

  % Forward radius inequality in the proof of Theorem 6.14.
  \item Prove absolute convergence of the differentiated series at
    \(|x-c|<R\) by choosing \(|x-c|<\rho<R\).
    Explain why the factors \(nr^{n-1}\), for \(0<r<1\),
    remain bounded and how this controls the extra factor of \(n\).
    Handle \(x=c\) separately.

  % Converse radius inequality, repairing the absolute-divergence gap.
  \item Prove the reverse radius inequality in Theorem 6.14:
    if the differentiated series converges at distance \(d>0\)
    from its center, show that the original series converges
    absolutely at every smaller distance. Why would divergence of
    an absolute-value series alone be insufficient to exclude
    conditional convergence of its signed series?

  % Theorem 6.15 and its proof.
  \item State and prove the theorem permitting differentiation of
    a convergent power series inside its radius (Theorem 6.15).
    Apply the differentiation theorem for function sequences to
    partial sums on smaller intervals. Why is uniform convergence
    on the entire open convergence interval unnecessary?

  % Theorem 6.16 and its proof.
  \item State the smoothness and coefficient-recovery theorem
    (Theorem 6.16). Derive the general formula for the \(k\)th
    derivative series and evaluate it at the center to recover
    each \(a_k\). What radius do all these derivative series have?

  % Corollary 6.17; use the common center c in the equality hypothesis.
  \item State and prove uniqueness of power-series coefficients
    when two convergent power series with the same center agree
    on a neighborhood of that center (Corollary 6.17).
    Why is a positive radius of convergence important?
\end{enumerate}

\subsection{The exponential function}
\textit{Source: \S6.5, pp.~82--84.}

\begin{enumerate}[resume]
  % Power-series definition and differentiation at the start of Section 6.5.
  \item Define \(E(x)\) by the exponential power series, without
    presupposing properties of \(e^x\). Prove that \(E\) is smooth
    on \(\R\), that \(E'=E\), and that \(E(0)=1\).

  % Proposition 6.18.
  \item State the addition law for \(E\) (Proposition 6.18).
    What rearrangement of the product of two power series appears
    in its proof, and why is that rearrangement justified?

  % Proof of Proposition 6.18.
  \item Prove the addition law by collecting terms of total degree
    \(n\) in the product \(E(x)E(y)\) and using the binomial theorem.
    Why would assuming the usual exponential addition law make
    this proof circular?

  % Consequences following Proposition 6.18.
  \item Derive the reciprocal relation between \(E(-x)\) and \(E(x)\).
    Prove that \(E(x)>0\) for every real \(x\), including negative inputs.

  % Proposition 6.19 and its proof.
  \item State and prove that \(E(x)\) grows faster than every
    nonnegative integer power of \(x\) as \(x\to\infty\)
    (Proposition 6.19). Which single positive term of its power
    series gives a useful lower bound? Extend the limit to any
    polynomial divided by \(E(x)\).

  % Proposition 6.20.
  \item State the initial-value characterization of the exponential
    (Proposition 6.20). Which differential equation and initial
    condition determine it uniquely on \(\R\)?

  % Proof of Proposition 6.20; use the correct quotient-rule numerator.
  \item Prove Proposition 6.20 by differentiating \(f/E\).
    Explain why the quotient is defined everywhere, why its
    derivative vanishes, and how the initial condition fixes the constant.

  % Closing paragraph of Section 6.5; combine established properties.
  \item How can the real logarithm be defined from \(E\)?
    Use positivity, the derivative, the addition law, growth,
    and continuity to justify that \(E\) is a bijection from
    \(\R\) onto \((0,\infty)\).
\end{enumerate}

\subsection{Taylor's theorem and power series}
\textit{Source: \S6.6, pp.~84--89.}

\subsubsection*{Smooth functions and analytic functions}
\begin{enumerate}[resume]
  % Opening comparison between Theorems 4.41 and 6.16.
  \item Distinguish the finite-order statement of Taylor's theorem
    from representation of a function by an infinite power series.
    Which limiting variable changes when the Taylor order is fixed,
    and which changes when the input is fixed?

  % Definition 6.21 and the following discussion.
  \item Define \(C^k(a,b)\) and \(C^\infty(a,b)\), including the
    meaning of smoothness. Why does existence of derivatives of
    every order force all of them to be continuous, while existence
    of only \(k\) derivatives need not make the \(k\)th derivative continuous?

  % Definition 6.22.
  \item Define a real analytic function on an interval.
    What is required at every center, and must one power series
    represent the function throughout the entire interval?

  % Analyticity implies smoothness, following Definition 6.22.
  \item Prove that every analytic function is smooth, and explain
    why its local power-series coefficients must be its Taylor
    coefficients. Does the converse follow from this argument?

  % Two failure modes discussed before Section 6.6.2.
  \item For a smooth function, define its formal Taylor series
    at a point. Describe the two possible failures discussed in
    the notes that prevent this formal series from representing
    the function in a neighborhood of the point.
\end{enumerate}

\subsubsection*{A smooth, non-analytic function}
\begin{enumerate}[resume]
  % Proposition 6.23; start with the basic difference quotient.
  \item Define
    \[
      \phi(x)=
      \begin{cases}
        \exp(-1/x),&x>0,\\
        0,&x\leq0.
      \end{cases}
    \]
    Prove continuity at zero and compute \(\phi'(0)\) from both
    sides. How does exponential growth at infinity control the
    right-hand difference quotient?

  % Induction in the proof of Proposition 6.23.
  \item Show inductively that, for \(x>0\),
    \(\phi^{(n)}(x)=p_n(1/x)\exp(-1/x)\) for a polynomial \(p_n\).
    Use this form to prove existence of every derivative at zero
    and determine all the values \(\phi^{(n)}(0)\).

  % Corollary 6.24.
  \item Find the Taylor series of \(\phi\) at zero, its radius of
    convergence, and its sum. Why does this establish that a smooth
    function need not be analytic, even when its Taylor series converges
    everywhere?

  % Discussion following Corollary 6.24.
  \item Why do the vanishing Taylor polynomials of \(\phi\) not
    contradict Taylor's theorem? Explain how bounds of the form
    \(0\leq\phi(x)\leq C_{n+1}x^{n+1}\), with constants depending
    on \(n\), describe behavior as \(x\to0^+\) without forcing
    \(\phi(x)=0\) at every fixed positive input.

  % Example 6.25.
  \item Give a smooth function that is positive for every nonzero
    real input and has all derivatives zero at zero.
    Construct it by composing \(\phi\) with \(x^2\), and explain
    its smoothness and failure of analyticity at zero.

  % Definition 6.26 and the piecewise-linear example on p. 88.
  \item Define compact support for a function on \(\R\).
    Give a nonzero continuous function with compact support and
    explain why a piecewise-linear triangular example is not smooth.

  % Example 6.27.
  \item Construct a nonzero smooth function with compact support
    using \(\phi(1-x^2)\). Write its piecewise formula, determine
    where it is positive and where it vanishes, and justify smoothness
    at the two transition points.

  % Closing discussion of Section 6.6; uniqueness is stated without proof.
  \item Compare what the value and all derivatives at one point
    determine for smooth functions and for analytic functions on
    an interval, using the uniqueness property quoted in the notes.
    Why can a nonzero analytic function on \(\R\) not have compact
    support? How does this apply to the smooth bump function above?
\end{enumerate}

\subsection{Appendix: Review of series}
\textit{Source: \S6.7, pp.~89--92.}

\begin{enumerate}[resume]
  % Definition 6.28.
  \item Define convergence and absolute convergence of a numerical
    series using partial sums. Distinguish a term, a partial sum,
    and the sum of the series.

  % Theorem 6.29 and its proof.
  \item State and prove the Cauchy criterion for a real series
    (Theorem 6.29). Specify which finite tails must satisfy the bound
    and explain the role of completeness of \(\R\).

  % Observation following Theorem 6.29.
  \item Prove that an absolutely convergent real series converges,
    using the Cauchy criterion and the triangle inequality.
    Give a convergent series that is not absolutely convergent.

  % Theorem 6.30.
  \item State and prove the necessary condition on the terms of
    a convergent series (Theorem 6.30). Is the condition sufficient?
    Give a counterexample to any invalid converse.

  % Theorem 6.31 and its proof.
  \item State and prove the comparison test (Theorem 6.31), including
    its sign and absolute-value conditions. Which series is shown
    to converge absolutely?

  % Theorem 6.32.
  \item State the ratio test (Theorem 6.32), including eventual
    nonvanishing of terms, existence of the absolute ratio limit,
    and the case where it diverges to infinity.
    What happens below, above, and at the threshold one?

  % Proof of Theorem 6.32.
  \item Prove the two decisive cases of the ratio test by comparison
    with geometric behavior. When the ratio limit is less than one,
    how is a geometric upper bound obtained? When it is greater
    than one, why do the terms fail to approach zero?

  % Definition 6.33 and the accompanying conventions.
  \item Define the limit superior of a real sequence by tail
    suprema. Why is the tail-supremum sequence decreasing?
    Explain the conventions giving finite, \(+\infty\), and
    \(-\infty\) values. Can a limit superior exist when the
    original sequence has no limit?

  % Proposition 6.34.
  \item State the eventual-upper-bound and infinitely-often-lower-bound
    characterization of a finite limit superior (Proposition 6.34).
    State also the versions for \(+\infty\) and \(-\infty\).
    Distinguish an eventual condition from one holding infinitely often.

  % Theorem 6.35.
  \item State the root test (Theorem 6.35) using
    \(\limsup_{n\to\infty}|a_n|^{1/n}\).
    Give the conclusions when this number is below one or above
    one, including the zero and infinite cases. What does the test
    say when it equals one?

  % Convergent case in the proof of Theorem 6.35.
  \item Prove the convergent case of the root test by choosing a
    number strictly between the limit superior and one.
    How does Proposition 6.34 yield a geometric bound on all
    sufficiently late terms?

  % Divergent case in the proof of Theorem 6.35.
  \item Prove the divergent case of the root test using a number
    strictly between one and the limit superior.
    Why are infinitely many large terms enough to establish
    divergence, even without a lower bound on every late term?

  % Synthesis of the appendix and its use in Sections 6.2--6.3.
  \item Compare the hypotheses of the ratio and root tests.
    Which accommodates zero terms and oscillating root magnitudes?
    Use the lacunary power series of \S6.3 to illustrate why the
    more general test is useful.
\end{enumerate}
